When a neural network overfits to an image, its thousands of weights suddenly become the image itself.
Store the weights, and you recover the picture.
This is the core claim of COIN~\citep{chen2022coin} and its successors, which treat neural network weights as a new medium for representing signals.
Yet a fundamental question remains unanswered: \emph{why} does this representation compress so well?
What geometric property makes the weight space of an image-fitted network amenable to low-dimensional description?

This paper provides the answer.
We study how the optimal parameters $\theta^*$ of an implicit neural representation (INR) move through weight space when the target signal undergoes group transformations.
Scale an image by factors from $0.5$ to $1.5$, or rotate it from $0^\circ$ to $360^\circ$, and refit the network each time.
The resulting parameter trajectories are not random wanderings in high-dimensional space.
Instead, they trace low-dimensional manifolds---orbits---whose geometry is dictated by the underlying transformation group.

Our first contribution is a tangent space analysis of these orbits.
For the rotation group $\SO(2)$, trajectories form closed ellipses; the parameters return to their starting point after a $360^\circ$ rotation.
For the scale group $\mathbb{R}^+$, trajectories extend monotonically along exponential paths and do not close.
We prove that the tangent vectors $\{J_X(f)\}$ form an approximate linear representation of the governing Lie algebra, and we quantify the deviation via a symmetry error $\delta$ with concrete empirical bounds.

Our second contribution is a practical application.
Standard meta-learning uses gradient directions as parameter update proxies, but gradients point orthogonal to the orbit---away from the manifold.
We derive a tangent-aware initialization using Jacobian-Vector Products that stays on the orbital manifold, reducing inner-loop convergence steps by $\ge 30\%$ in conditional INR adaptation while maintaining final reconstruction quality.

By connecting the vertical geometry of parameter orbits to the empirical success of COIN++ modulation-based compression, this work lays the geometric foundation for understanding and exploiting the structure of weight space in implicit neural representations.
